\scanpage{122}
% archive sheet number: 121

$\mathbb{X}$ \deleted{et}\inserted{\handwritten{(}donc} $X$ est régulier en les points de son support\handwritten{)},\deleted{$U=f^{\handwritten{-1}}(V)\handwritten{\unsure{(}}=X-\operatorname{supp}D\handwritten{\unsure{)}}$}.
\margin{\handwritten{\uline{(Condition de transversalité)}-}}

\handwritten{\unsure{b|a}}) $\operatorname{supp}E$ est irréductible, et il est régulier en les points de $f(\operatorname{supp}D)=\operatorname{supp}E\handwritten{\cap}f(X)$ (donc $Y$ est régulier en ses points).

Alors pour tout $x\handwritten{\in}\operatorname{supp}D$, \deleted{\unsure{l'image des}}\inserted{les} sous-groupes d'inertie de $G$ en $x$ \handwritten{o}nt comme images par $\handwritten{\phi}:G\handwritten{\longrightarrow}H$ des sous-groupes de $H$ qui sont des sous-groupes d'inertie de $H$ relativement au point maximal de $\operatorname{supp}E$. Par suite, pour toute représentation linéaire de $G$ (dans un module $M$ sur quelque anneau) \deleted{les} qui veut bien se factoriser par une représentation linéaire \uline{de} $H$, les images dans $\operatorname{Aut}(M)$ des \inserted{sous-}groupes d'inertie \deleted{\unsure{dans}} de $G$ relatifs aux divers points de $\operatorname{supp}D$ (ou encore, relatifs aux diverses composantes irréductibles de $\operatorname{supp}D$, cela revient au même en vertu de 3\handwritten{)} sont conjugués dans $\operatorname{Aut}(M)$.

\inserted{\uline{Remords\handwritten{\unsure{:|;}} au n\textsuperscript{o}}\handwritten{s3,4}.} On peut demander \deleted{\unsure{un}}\handwritten{\unsure{des|du}} résultat\handwritten{\unsure{s}} plus précis, en \deleted{\uline{Application}} considérant pour tout $\handwritten{\xi}$ comme ci-dessus (au dessus d'un point géométrique\inserted{$\handwritten{\bar}{x}$} de $\operatorname{supp}D$) l'homomorphisme correspondant $\underline{Z}_{\handwritten{\ell}}(1)(\handwritten{\bar}{x})\handwritten{\longrightarrow}G$ dont l'image est le sous-groupe d'inertie correspondant, et en se demandant comment cet homomorphisme varie avec le choix de $\handwritten{\xi}$ \uline{et de} $\handwritten{\bar}{x}$. P\unsure{$\overset{\smile}{\underset{\frown}{\mathrm{o}}}$}ur $\handwritten{\xi}$ variable, la variation se fait évidemment modulo composition avec les automorphismes intérieurs de $G$, \deleted{Pour donner un sens à la question} donc pour $\handwritten{\bar}{x}$ fixé, on trouve une classe de conjugaison d'homomorphismes
\[
\underline{Z}_{\handwritten{\ell}}(1)(\handwritten{\bar}{x})\handwritten{\longrightarrow}G\ .
\]
\deleted{pour}\inserted{\handwritten{\unsure{Étudions|Étudier}}} \deleted{L}\handwritten{l}a variation de cet homomorphisme\handwritten{,} pour $\handwritten{\bar}{x}$ variable dans une composante \deleted{de $D$}\inserted{$D_i$ \handwritten{de}} l'ensemble des points réguliers de $D$\inserted{à car. $\handwritten{\ne\ell}$}; d'après ce qu'on a vu au n\textsuperscript{\uline{o}}\deleted{\unsure{2}}\handwritten{3}, on passe d'un tel homomorphisme à un autre en composant avec un quelconque des isomorphismes $\underline{Z}_{\handwritten{\ell}}(1)(\handwritten{\bar}{x})\handwritten{\longrightarrow}\underline{\unsure{\mathbb{Z}|\deleted{E}Z|\deleted{Z}E|E|Z}}_{\handwritten{\ell}}(1)(\handwritten{\bar}{x}')$ associé à une classe de chemins de $\handwritten{\bar}{x}$ à $\handwritten{\overline}{x'}$. Si cet isomorphe ne dépend pas du choix de cette classe de chemins, i.e. si l\unsure{\deleted{'}e|é}\deleted{action du groupe fondamental} faisceau l-adique de \deleted{\unsure{t}}\handwritten{T}ate $\underline{Z}_{\handwritten{\ell}}(1)$ sur $D(\handwritten{\ell})$ est trivial (ce qui se voit déjà sur la

\figure{\handwritten{\begin{tikzpicture}[x=0.006cm,y=-0.006cm]
\draw (1500,0) .. controls (1380,0) and (1340,24) .. (1200,18) .. controls (1000,12) and (810,8) .. (655,16) .. controls (625,20) and (605,18) .. (565,15);
\draw (553,16) -- (506,11) -- (484,90) .. controls (443,86) and (416,94) .. (377,92) .. controls (310,99) and (260,119) .. (222,119) -- (147,107) -- (16,120);
\draw (16,120) -- (24,116) (16,120) -- (24,124);
\draw (377,92) -- (385,88) (377,92) -- (385,96);
\draw (14,120) .. controls (10,174) and (23,234) .. (15,288) .. controls (19,365) and (7,421) .. (13,492) .. controls (10,568) and (17,612) .. (17,663) .. controls (19,728) and (20,790) .. (29,851) .. controls (28,914) and (50,953) .. (40,1011) .. controls (25,1050) and (28,1089) .. (34,1142) -- (39,1186);
\draw[line width=0.3pt] (25,138) .. controls (15,199) and (19,253) .. (22,316) .. controls (30,386) and (12,452) .. (17,486);
\draw[line width=0.3pt] (3,623) .. controls (0,655) and (3,676) .. (2,701) .. controls (-1,726) and (4,749) .. (5,774) .. controls (10,801) and (8,829) .. (11,852) .. controls (13,878) and (19,895) .. (22,917) .. controls (23,946) and (25,976) .. (23,994) .. controls (21,1012) and (27,1032) .. (27,1055) .. controls (26,1076) and (32,1096) .. (29,1118) .. controls (29,1141) and (33,1160) .. (33,1180) .. controls (34,1193) and (42,1208) .. (44,1221);
\end{tikzpicture}}}
